\begin{tikzpicture}[font=\small,line cap=round,line join=round]\draw[otRule!65,line width=.2pt](0,0.000)--(7.10,0.000);\node[left,font=\tiny]at(0,0.000){$0$};\draw[otRule!65,line width=.2pt](0,0.680)--(7.10,0.680);\node[left,font=\tiny]at(0,0.680){$10$};\draw[otRule!65,line width=.2pt](0,1.360)--(7.10,1.360);\node[left,font=\tiny]at(0,1.360){$20$};\draw[otRule!65,line width=.2pt](0,2.040)--(7.10,2.040);\node[left,font=\tiny]at(0,2.040){$30$};\draw[otRule!65,line width=.2pt](0,2.720)--(7.10,2.720);\node[left,font=\tiny]at(0,2.720){$40$};\draw[otRule!65,line width=.2pt](0,3.400)--(7.10,3.400);\node[left,font=\tiny]at(0,3.400){$50$};\draw[otMuted,->](0,0)--(0,3.80)node[above,font=\scriptsize]{piles};\draw[otMuted](0,0)--(7.10,0);\node[below,font=\scriptsize]at(1.00,0){groupe 1};\fill[otNavy!60](0.790,0)rectangle(1.210,1.564);\node[above,font=\tiny,text=otInk]at(1.000,1.564){23};\geoSolution{\node[above,font=\scriptsize,text=otGreenDark]at(1.000,1.864){$0{,}46$};}\node[below,font=\scriptsize]at(2.70,0){groupe 2};\fill[otNavy!60](2.490,0)rectangle(2.910,1.904);\node[above,font=\tiny,text=otInk]at(2.700,1.904){28};\geoSolution{\node[above,font=\scriptsize,text=otGreenDark]at(2.700,2.204){$0{,}56$};}\node[below,font=\scriptsize]at(4.40,0){groupe 3};\fill[otNavy!60](4.190,0)rectangle(4.610,1.700);\node[above,font=\tiny,text=otInk]at(4.400,1.700){25};\geoSolution{\node[above,font=\scriptsize,text=otGreenDark]at(4.400,2.000){$0{,}50$};}\node[below,font=\scriptsize]at(6.10,0){groupe 4};\fill[otNavy!60](5.890,0)rectangle(6.310,2.108);\node[above,font=\tiny,text=otInk]at(6.100,2.108){31};\geoSolution{\node[above,font=\scriptsize,text=otGreenDark]at(6.100,2.408){$0{,}62$};}\draw[otOrange,dashed,line width=.8pt](0,1.700)--(7.10,1.700)node[right,font=\scriptsize,text=otInk]{25 : moitié};\fill[otNavy!60](7.40,3.20)rectangle++(.3,.25);\node[anchor=west,font=\scriptsize]at(7.75,3.32){piles sur 50 lancers};\end{tikzpicture}\par\vspace{1mm}{\fontsize{8}{10}\selectfont\color{otMuted}Nombre de piles obtenus par chaque groupe ; au corrigé, les fréquences sont indiquées en vert.}